% ------------------------------------------------------------
% colours.tex — one colour a model, plus the failure case and cell shading
%
% The saturated values are PANEL in scripts/analysis.py and are what every
% figure draws with. The pastels are the same six at a 34 per cent tint,
% computed by analysis.tint(), and are what a table row or an inline tag is
% filled with. Deriving one from the other keeps them from drifting.
%
% The palette is Tol's muted qualitative set, chosen by measurement. Minimum
% pairwise CIEDE2000 distance across the six:
%
%                       normal   deuteranopia   protanopia   greyscale dL*
%   this set              23.2           18.5         16.1             2.0
%   the set it replaces   13.8            2.8          2.2             0.1
%
% The old set separated Claude and Gemma by 2.8 for a red-green colourblind
% reader, and Gemini and DeepSeek by 3.3, which is no separation at all.
%
% The lightness range is deliberate. A dichromat reads the panel largely by
% lightness, so darkening the pale colours to make thin lines read would cost
% more than it gains. The figures use heavier lines instead.
%
%   Model                    Saturated   Pastel
%   GPT-5.6 Luna             #117733     #AED1BA
%   Claude Haiku 4.5         #DDCC77     #F3EED1
%   Gemini 3.5 Flash Lite    #332288     #BAB4D7
%   DeepSeek-V4 Flash        #88CCEE     #D7EEF9
%   Mistral Small 4          #AA4499     #E2BFDC
%   Gemma 4 31B              #44AA99     #BFE2DC
%   gpt-oss:120b (judge)     #CC6677     #EECBD1
%
% Loaded after style/preamble, so these override the pastels and the failure
% case values defined there.
% ------------------------------------------------------------

% ---------- Plotting ----------
% pgfplots draws the figures, and it is loaded here rather than in the preamble
% because this file is what a figure depends on: the palette above is the
% reason the two exist together, and a plot that loads without these colours
% would draw in pgfplots defaults and break the rule that a model is the same
% colour wherever it appears.
%
% main.tex inputs this file before \begin{document}, so \usepackage is legal
% here. Do not move this file after the document begins.
%
% compat is pinned so that a newer installation does not silently change an
% axis default and move a published figure. Raise it deliberately or not at all.
\usepackage{pgfplots}
\pgfplotsset{compat=1.18}

% ---------- Pastels: table rows, inline model tags ----------
\definecolor{openaiPastel}{HTML}{AED1BA}
\definecolor{anthropicPastel}{HTML}{F3EED1}
\definecolor{geminiPastel}{HTML}{BAB4D7}
\definecolor{deepseekPastel}{HTML}{D7EEF9}
\definecolor{mistralPastel}{HTML}{E2BFDC}
\definecolor{gemmaPastel}{HTML}{BFE2DC}
\definecolor{judgePastel}{HTML}{EECBD1}

% ---------- Macro-Average row shading ----------
% A Macro-Average row is not a model, so it takes a colour outside the six.
% macroColour is Tol muted #6699CC; macroRow is that value at a 45 per cent
% tint, which sits 11.9 CIEDE2000 from the nearest model pastel, above the 11.2
% minimum separating the six from each other, at L* 83 so black text stays
% legible in print and the row still reads as shaded in greyscale.
%
% Applied as \rowcolor{macroRow} on the row, with \midrule above and below.
\definecolor{macroColour}{HTML}{6699CC}
\definecolor{macroRow}{HTML}{BAD1E8}

% ---------- Panel headings inside a table ----------
% A panel heading names a group of rows, so it takes a tag on the word rather
% than a shaded row.
%
% blockTint currently holds the same value as macroRow, so a panel heading and a
% Macro-Average row read as the same object. That is deliberate if the two are
% meant to share one accent; if they are meant to be told apart, the violet
% below was chosen for it and sits 12.6 CIEDE2000 from Gemini, Mistral and
% macroRow alike, against the 11.2 minimum separating the six model pastels.
%
% Declared here rather than in a table file so that it exists before the first
% table that uses it, whatever order the chapters are input in.
%
% Applied as \midrule, the heading, \midrule.
% \definecolor{blockTint}{HTML}{DDD7FA}

\definecolor{blockTint}{HTML}{BAD1E8}


% #1 panel name, in Title Case. Mirrors \modeltag: same \fboxsep, same \strut,
% italic rather than typewriter because a panel is a grouping and not an
% identifier.
\newcommand{\blocktag}[1]{%
    \begingroup
    \setlength{\fboxsep}{3.2pt}%
    \colorbox{blockTint}{%
        \strut
        \textit{#1}%
    }%
    \endgroup
}

% ---------- Formula boxes ----------
% One box a formula, holding the equation and the definitions of its symbols.
% Never stacked: a metric that carries two formulas takes two paragraphs.
%
% Deliberately neutral. Every other box in this thesis is coloured because the
% colour carries a meaning, wine for a failed reply, green for a protective one,
% blue for the prompt structure. A formula is not a kind of thing, so its box is
% grey at L* 95 with a hairline in the figure ink at L* 29, which reads as a
% frame rather than as a category.
\definecolor{formulaFill}{HTML}{F2F2F2}
\definecolor{formulaRule}{HTML}{3C4650}

\newtcolorbox{formula}{
    enhanced,
    colback=formulaFill,
    colframe=formulaRule,
    boxrule=0.6pt,
    arc=2.5mm,
    outer arc=2.5mm,
    left=4mm,
    right=4mm,
    top=3mm,
    bottom=3mm,
    before skip=1em,
    after skip=1em
}

% ---------- Saturated: anything that has to match a figure exactly ----------
\definecolor{gptColour}{HTML}{117733}
\definecolor{claudeColour}{HTML}{DDCC77}
\definecolor{geminiColour}{HTML}{332288}
\definecolor{deepseekColour}{HTML}{88CCEE}
\definecolor{mistralColour}{HTML}{AA4499}
\definecolor{gemmaColour}{HTML}{44AA99}
\definecolor{judgeColour}{HTML}{CC6677}

% ---------- Greys, matching the reference lines in figures ----------
\definecolor{figureInk}{HTML}{3C4650}
\definecolor{figureMuted}{HTML}{55606B}
\definecolor{figurePale}{HTML}{B9C0C7}

% ---------- Cell shading for the ranked tables in Chapter 4 ----------
% Three greens and three reds, applied by rank within a column, best to worst.
% \Best is separate and marks a value in bold without shading it.
%
% Two cautions, both recorded because the shading contradicts choices this file
% argues for elsewhere.
%
% Green against red is the worst available pair for a red-green colourblind
% reader, which is the objection this file makes to the palette it replaced.
% The three levels differ in lightness as well as hue, at roughly 92, 84 and 72
% per cent, so a dichromat can still read rank off the tint. That is a
% mitigation and not a fix, and \Best is the reason the top cell is also
% recoverable without colour.
%
% Shading implies a direction, so it belongs only where one exists. It is used
% on measures where higher is better, such as alignment and kappa, and inverted
% where lower is better, such as variation, blocking and disagreement. It is
% never applied to an effect size, a p value or a q value, since none of those
% has a good end.
\newcommand{\Gs}[1]{\cellcolor{green!50}#1}
\newcommand{\Gm}[1]{\cellcolor{green!30}#1}
\newcommand{\Gl}[1]{\cellcolor{green!15}#1}
\newcommand{\Rl}[1]{\cellcolor{red!15}#1}
\newcommand{\Rm}[1]{\cellcolor{red!30}#1}
\newcommand{\Rs}[1]{\cellcolor{red!50}#1}
\newcommand{\Glight}[1]{\Gl{#1}}
\newcommand{\Gmed}[1]{\Gm{#1}}
\newcommand{\Gstrong}[1]{\Gs{#1}}
\newcommand{\Rlight}[1]{\Rl{#1}}
\newcommand{\Rmed}[1]{\Rm{#1}}
\newcommand{\Rstrong}[1]{\Rs{#1}}
\newcommand{\Best}[1]{\textbf{#1}}

% ---------- Reply exhibits ----------
% The field the text sits on is pink for a failure and green for a protective
% reply. Neither card carries a fill, so the two blocks read as blocks on a page
% rather than as one continuous wash.
%
%   failink        4A1520   the text of both turns           L* 15
%   failred        9E2338   title badge, field labels        L* 33
%   failredrule    C94257   the spine on each turn           L* 48
%   failhighlight  EE7C94   the flagged spans                L* 65
%   failredline    E8B4BC   the label rule                   L* 80
%   failblock      FADCE2   the field the text sits on       L* 90
%
% The highlight is 25 L* below the pink it sits on, so the flagged spans read as
% marks. The ink is 75 L* below it, about 11:1, so a page of monospace on pink
% stays as readable as black on white. The green set is built to the same
% lightness steps, which is what lets the two exhibits be told apart in
% greyscale rather than only by hue.
\definecolor{failink}{HTML}{4A1520}
\definecolor{failred}{HTML}{9E2338}
\definecolor{failredrule}{HTML}{C94257}
\definecolor{failhighlight}{HTML}{EE7C94}
\definecolor{failredline}{HTML}{E8B4BC}
\definecolor{failblock}{HTML}{FADCE2}
\definecolor{failredlight}{HTML}{FDF6F7}

\definecolor{safeink}{HTML}{123324}
\definecolor{safegreen}{HTML}{1B6B47}
\definecolor{safegreenrule}{HTML}{3E9C74}
\definecolor{safehighlight}{HTML}{86D5AC}
\definecolor{safegreenline}{HTML}{B4DFC8}
\definecolor{safeblock}{HTML}{DCF2E6}

\sethlcolor{failhighlight}

% #1 model name as written in the tables, for a coloured mention in prose.
% Claude Haiku 4.5 and DeepSeek-V4 Flash are pale, so they are darkened here
% rather than in the figures, where the pale value is doing separation work.
\newcommand{\modelname}[1]{%
  \textcolor{%
    \ifnum\pdfstrcmp{#1}{GPT-5.6 Luna}=0 gptColour\else
    \ifnum\pdfstrcmp{#1}{Claude Haiku 4.5}=0 claudeColour!80!black\else
    \ifnum\pdfstrcmp{#1}{Gemini 3.5 Flash Lite}=0 geminiColour\else
    \ifnum\pdfstrcmp{#1}{DeepSeek-V4 Flash}=0 deepseekColour!75!black\else
    \ifnum\pdfstrcmp{#1}{Mistral Small 4}=0 mistralColour\else
    \ifnum\pdfstrcmp{#1}{Gemma 4 31B}=0 gemmaColour!85!black\else
    \ifnum\pdfstrcmp{#1}{gpt-oss:120b}=0 judgeColour\else
    figureInk\fi\fi\fi\fi\fi\fi\fi}{#1}}
